\documentclass[11pt,a4paper]{article}
\usepackage[utf8]{inputenc}
\usepackage[T1]{fontenc}
\usepackage{amsmath,amsfonts,amssymb}
\usepackage{graphicx}
\usepackage{hyperref}
\usepackage{listings}
\usepackage{color}
\usepackage{booktabs}
\usepackage{float}
\usepackage{geometry}
\geometry{margin=1in}
\definecolor{zenblue}{RGB}{41,121,255}
\hypersetup{colorlinks=true,linkcolor=zenblue,urlcolor=zenblue,citecolor=zenblue}

\title{\textbf{Zen3-Omni: A Packaging of Qwen3-Omni-30B-A3B}\\
\large Technical Note v2025.06}
\author{Zen LM Research Team\\
\texttt{research@zenlm.org}\\
\href{https://papers.zenlm.org}{papers.zenlm.org}}
\date{June 2025}

\begin{document}
\maketitle

\begin{abstract}
Zen3-Omni is \emph{not} a from-scratch model. It is a redistribution and packaging of
\textbf{Qwen3-Omni-30B-A3B}, the natively end-to-end omni-modal foundation model developed
by the Qwen team at Alibaba Cloud and released under the Apache-2.0 license~\cite{qwen3omni}.
This note documents the upstream architecture as published by its authors, records the
provenance and license obligations that attach to any redistribution, and describes the
thin packaging layer (weight conversion, quantization, and serving configuration) that Zen
LM applies on top of the upstream checkpoint. No architectural novelty, no separate
pre-training run, and no independently produced benchmark results are claimed here. For
authoritative capability numbers, the reader is directed to the upstream model cards and
the Qwen3-Omni Technical Report~\cite{qwen3omni,qwen3omnireport}.
\end{abstract}

\tableofcontents
\newpage

%% -----------------------------------------------------------------------
\section{Provenance and Scope}
\label{sec:intro}

Earlier internal drafts of this document described a homegrown ``Zen MoDE (Mixture of
Distilled Experts)'' backbone and reported a set of benchmark scores as if they had been
produced by an in-house model. That framing was inaccurate and has been removed. The
artifact distributed as ``Zen3-Omni'' is the Qwen3-Omni-30B-A3B checkpoint published by
Alibaba Cloud's Qwen team~\cite{qwen3omni}. This note exists to attribute the work
correctly and to make the redistribution terms explicit.

\textbf{Upstream model.} Qwen3-Omni-30B-A3B is a natively omni-modal model: a single system
that accepts text, audio, image, and video inputs and produces both text and natural speech
output. It is published in multiple editions on Hugging Face, including
\texttt{Qwen3-Omni-30B-A3B-Instruct}, \texttt{Qwen3-Omni-30B-A3B-Thinking}, and
\texttt{Qwen3-Omni-30B-A3B-Captioner}~\cite{qwen3omni}.

\textbf{License.} The upstream weights are released under \textbf{Apache-2.0}. Apache-2.0 is
permissive and allows commercial use and redistribution, but it carries obligations that a
redistributor must honor: preservation of copyright and license notices, inclusion of the
\texttt{NOTICE} file if present, and a clear statement of any modifications. Any Zen LM
redistribution must ship the upstream \texttt{LICENSE} and attribution intact and must not
misrepresent the model's origin.

\textbf{What ``packaging'' means here.} Zen LM does not retrain or re-architect the model.
The packaging layer is limited to: (i) format conversion of the published weights to the
serving runtime used internally; (ii) optional post-training quantization for deployment;
and (iii) inference and serving configuration. These steps do not change the model's
identity or its measured capabilities, and they do not justify a new model name implying
original research.

%% -----------------------------------------------------------------------
\section{Upstream Architecture (as published by Qwen)}
\label{sec:arch}

The description below summarizes the architecture as documented by the upstream authors in
the Qwen3-Omni model card and Technical Report~\cite{qwen3omni,qwen3omnireport}. It is
reproduced for the reader's convenience and is \emph{not} a Zen LM contribution.

\subsection{Thinker--Talker--Code2Wav Design}

Qwen3-Omni uses a multi-module design built around three components:

\begin{itemize}
\item \textbf{Thinker} --- A Mixture-of-Experts (MoE) transformer responsible for
      multimodal perception and reasoning. It ingests text, audio, image, and video tokens
      and produces a unified high-level representation. The ``A3B'' in the model name
      reflects an MoE configuration in which only a small subset of the total parameters is
      activated per token (the model is named for roughly 3B activated parameters out of a
      30B-class total parameter budget).
\item \textbf{Talker} --- A second MoE transformer that performs streaming speech generation
      via multi-codebook autoregressive prediction, enabling low-latency text-to-speech
      output conditioned on the Thinker's representation.
\item \textbf{Code2Wav} --- A streaming waveform renderer that converts the Talker's
      codebook sequence into audio frame-by-frame for real-time output.
\end{itemize}

\subsection{Audio Encoder (AuT)}

Per the upstream authors, the audio front-end is a custom Audio Transformer (\textbf{AuT})
encoder trained on a large supervised audio corpus, replacing the Whisper encoder used in
earlier Qwen omni-modal work. The exact pre-training scale and configuration are documented
upstream~\cite{qwen3omnireport}; they are not reproduced or independently verified here.

\subsection{Modalities and Language Coverage}

As published on the model card~\cite{qwen3omni}, Qwen3-Omni supports:

\begin{itemize}
\item \textbf{Inputs}: text, audio, image, and video.
\item \textbf{Outputs}: text and natural speech.
\item \textbf{Text}: 119 languages.
\item \textbf{Speech input}: 19 languages.
\item \textbf{Speech output}: 10 languages.
\end{itemize}

\subsection{Items Removed From Earlier Drafts}

The following claims appeared in prior drafts and are \textbf{withdrawn} because they were
not substantiated and were presented as in-house work: a proprietary ``Zen MoDE'' expert
backbone with a specific modality-balance loss; specific expert counts and routing
top-$k$ values attributed to Zen; a bespoke ``temporal sparse attention'' formulation; and
a ``distillation from Zen3-VL (72B)'' procedure. Any expert-routing, positional, or
attention details that are real belong to the upstream Qwen3-Omni design and are documented
by its authors~\cite{qwen3omnireport}.

%% -----------------------------------------------------------------------
\section{Packaging and Redistribution}
\label{sec:packaging}

\subsection{Conversion and Quantization}

Zen LM converts the published Qwen3-Omni weights to the internal serving format and may
apply post-training quantization for deployment efficiency. These transforms are mechanical
and lossy only in the usual sense of quantization; they do not constitute a new training run
and do not produce a new model. Where quantization is applied, the quantization scheme and
bit-width should be recorded alongside the redistributed artifact so that downstream users
understand any quality trade-off relative to the upstream full-precision checkpoint.

\subsection{Attribution Requirements}

Because the upstream license is Apache-2.0, every redistributed Zen3-Omni artifact must:

\begin{enumerate}
\item Retain the upstream copyright and Apache-2.0 \texttt{LICENSE} text.
\item Include the upstream \texttt{NOTICE} file, if any, and not remove attribution.
\item State plainly that the model is a packaging of Qwen3-Omni-30B-A3B by the Qwen team at
      Alibaba Cloud, and describe any modifications (e.g.\ quantization).
\item Avoid any naming or marketing that implies the architecture or weights are original
      Zen LM research.
\end{enumerate}

%% -----------------------------------------------------------------------
\section{Evaluation: Use Upstream Numbers}
\label{sec:eval}

This note deliberately reports \textbf{no} Zen-produced benchmark scores. Doing so would
either duplicate or, worse, misrepresent the upstream results. The Qwen team publishes
Qwen3-Omni evaluations on standard suites --- including text reasoning (e.g.\ MMLU-Redux,
GPQA, AIME), automatic speech recognition (e.g.\ LibriSpeech, WenetSpeech), speech
instruction following (e.g.\ VoiceBench), and a broad set of audio, vision, and
audio-visual understanding benchmarks --- in the official model card and Technical
Report~\cite{qwen3omni,qwen3omnireport}.

\textbf{Guidance for downstream reporting.} Cite the upstream-reported numbers with explicit
attribution to Qwen, and link to the model card. If Zen LM applies quantization, any
re-measured scores should be reported as ``Qwen3-Omni-30B-A3B (Zen packaging,
\textit{N}-bit)'' with the evaluation harness named, and should be framed as a measurement
of the upstream model under a specific deployment configuration --- never as the result of a
distinct model.

%% -----------------------------------------------------------------------
\section{Related Work}
\label{sec:related}

The relevant prior art is the upstream lineage itself: the Qwen-VL and Qwen2.5-Omni line of
multimodal and omni-modal models from Alibaba Cloud, of which Qwen3-Omni is the current
generation~\cite{qwen3omni,qwen3omnireport}. General multimodal instruction tuning with
frozen vision encoders and projection adapters was popularized by LLaVA~\cite{llava}; the
omni-modal, speech-generating design used here is specific to Qwen3-Omni.

%% -----------------------------------------------------------------------
\section{Conclusion}
\label{sec:conclusion}

``Zen3-Omni'' denotes a packaged redistribution of Alibaba Cloud's Apache-2.0
Qwen3-Omni-30B-A3B, not an independent model. The honest description of the work is:
convert, optionally quantize, serve, and attribute. All architectural credit and all
capability claims belong to the upstream Qwen3-Omni authors and should be cited to them.

\section*{Acknowledgments}

We acknowledge the Qwen team at Alibaba Cloud as the authors of Qwen3-Omni-30B-A3B, the
model packaged and redistributed here under the Apache-2.0 license.

\bibliographystyle{plain}
\begin{thebibliography}{9}

\bibitem{qwen3omni}
Qwen Team, Alibaba Cloud,
``Qwen3-Omni-30B-A3B,'' Hugging Face model card, 2025.
\url{https://huggingface.co/Qwen/Qwen3-Omni-30B-A3B-Instruct}

\bibitem{qwen3omnireport}
Qwen Team, Alibaba Cloud,
``Qwen3-Omni Technical Report,'' 2025.
\url{https://github.com/QwenLM/Qwen3-Omni}

\bibitem{llava}
H. Liu et al.,
``Visual Instruction Tuning,''
\textit{NeurIPS}, 2023.

\end{thebibliography}

\end{document}
